\scanpage{365}
\mspage{104}

$C\times_T T'$, ou par $C_{/t_0}\simeq C\times_T T_{/t_0}$. \deleted{\unsure{Dans ce cas}} \inserted{Bien} mieux, elle ne changerait pas en remplaçant $T$ par $(\mathbb N)^\circ$, la catégorie (ordonnée) opposée à $\mathbb N$ (ayant donc un unique objet final), et la catégorie fibrée $C$ sur $T$ par la catégorie $C_{\underline t}$ induite par $C$ sur $(\mathbb N)^\circ$, via le foncteur
\[
\underline t:(\mathbb N)^\circ\mathrel{\unsure{-|\to}}T
\]
défini par le système projectif $\underline t=(t_n)_{n\in\mathbb N}$ dans $T$. Cela signifierait donc que pour comprendre le type homotopique (au sens strict, i.e. dans $\mathrm{Cat\,hot}$, ou large i.e. dans $\mathrm{Hot}$) de $C$, on peut \unsure{le|la} ``nettoyer'' de façon draconienne, la base $T$ étant \inserted{vue comme} une sorte de gangue un peu \unsure{vaseuse|vasarde}, ou des \unsure{scories} \unsure{pléthoriques},

\textcircled{1} De toutes façons, si $T'$ est une \inserted{autre} catégorie filtrante à gauche, et
\[
\unsure{g|\varphi}:T'\longrightarrow T
\]
un foncteur cofinal à gauche, \deleted{\unsure{(de}} de sorte que $C\deleted{\unsure{/T}}$ sur $T$, et $C'=C\times_T T'$ sur $T'$, définissent \inserted{essent.} le ``$=$'' objet de $\mathrm{Ind\,Cat\,hot}$, et que \unsure{la catégorie|les catégories} $C'_\infty$ \inserted{pour $C'$ sur $T'$} \unsure{fait|font} partie des catégories $C_\infty$ pour $C$ sur $T$ -- on peut se demander

% archive sheet number: 364
